\section*{Описание архитектуры решения}

Программа состоит из двух отдельных исполняемых файлов: \texttt{parent.exe} и \texttt{child.exe}. Взаимодействие организовано через два анонимных канала (pipe).

\begin{enumerate}
    \item Родительский процесс создает два канала: \texttt{pipe1} (для передачи чисел) и \texttt{pipe2} (для получения результата).
    \item С помощью \texttt{SetHandleInformation()} отключается наследование концов каналов, которые остаются у родителя. Это предотвращает взаимоблокировку (deadlock), так как дочерний процесс не будет держать канал открытым.
    \item Создается дочерний процесс через \texttt{CreateProcess()}. В структуре \texttt{STARTUPINFO} стандартный ввод потомка связывается с \texttt{pipe1\_read}, а стандартный вывод — с \texttt{pipe2\_write}.
    \item Родитель запрашивает у пользователя строку с числами, записывает её в \texttt{pipe1} и закрывает канал. Закрытие канала служит сигналом EOF для потомка.
    \item Дочерний процесс считывает числа из \texttt{stdin}, вычисляет сумму, записывает её в указанный файл и отправляет результат родителю через \texttt{stdout}.
    \item Родитель читает результат из \texttt{pipe2}, выводит его на экран и ожидает завершения потомка через \texttt{WaitForSingleObject()}, предотвращая утечку ресурсов.
\end{enumerate}

\subsection*{Кроссплатформенность}
Код разделен на платформенно-зависимую и платформенно-независимую части. Файл \texttt{platform.h} содержит абстрактные типы и функции. Для Windows реализация находится в \texttt{platform\_win.c}, для Linux — в \texttt{platform\_linux.c}. Основная логика не зависит от ОС.